\section{Résultats}

\subsection{Expérience 1 : Analyse DMA}

À la fin de chaque essai pour une fréquence donnée, l’instrument NETZSCH fournit les données
brutes de l'analyse mécanique dynamique. La Fig.~\ref{fig:DMA_brut} présente un exemple typique
des signaux obtenus : le module de stockage $E'$ (en bleu), le module de perte $E''$ (en jaune)
ainsi que le facteur d’amortissement $\tan\delta$ (en vert), mesurés lors d’une montée en
température entre 70\,°C et 175\,°C.

\vspace{0.4cm}
\begin{flushleft}
\begin{minipage}[2cm]{0.3\textwidth}
Le protocole est répété pour plusieurs fréquences, comprises entre 1~Hz et 20~Hz, sur des
échantillons de PMMA en configuration de flexion trois points. Pour chaque fréquence, nous
obtenons les courbes de $E'$ et de $\tan\delta$ en fonction de la température. Ces résultats
comparatifs sont illustrés en Fig.~\ref{fig:DMA_resultats}.
\end{minipage}
\hfill
\begin{minipage}{0.65\textwidth}
    \includegraphics[width=\linewidth]{Seite_graph_primitif.png}
	\captionof{figure}{Exemple de données brutes issues de la mesure DMA}
    \label{fig:DMA_brut}
\end{minipage}
\end{flushleft}
\\
\begin{flushleft}
\begin{minipage}[3cm]{0.34\textwidth}
La détection automatique du pic de $\tan\delta$ permet d’identifier la position de la
transition vitreuse pour chaque fréquence. Ces températures caractéristiques sont ensuite
utilisées pour construire le graphe d’Arrhenius et déterminer l’énergie d’activation du PMMA.
\end{minipage}
\hfill
\begin{minipage}{0.62\textwidth}
    \centering
    \includegraphics[width=\linewidth]{Seite_Peak_analyse.pdf}
    \captionof{figure}{Analyse du pic de $\tan\delta$ permettant l’extraction de $T_g(f)$}
    \label{fig:DMA_resultats}
\end{minipage}
\end{flushleft}


\\
\newpage

\subsection{Expérience 2, DSC}

\begin{figure}[hct]
    \centering
	\includegraphics[scale=0.4]{fig1.png}
	\captionof{figure}{donees brutes}
        \label{fig:DSCprimitif}
\end{figure}


\begin{minipage}[t]{0.5\textwidth} 
	
\includegraphics[width=\textwidth]{dsc_chauffe.pdf}%
\captionof{figure}{donees brutes}
       \label{fig:DSCprimitif}
\end{minipage}
\hfill
\begin{minipage}[h]{0.4\textwidth}

BLBBLBAAAA AAAAAAA AA AAAAAAAAAA,
skfvnlsnvf,sljvj
;sfi;vn;s
	
\end{minipage}


\begin{minipage}[t]{0.45\textwidth}
    \centering
    \includegraphics[width=\linewidth]{dsc_amorphe.pdf}
        \captionof{figure}{donees brutes}
        \label{fig:DSCprimitif}
\end{minipage}
\hfill
\begin{minipage}[h]{0.4\textwidth}
BLBBLBAAAA AAAAAAA AA AAAAAAAAAA,
skfvnlsnvf,sljvj
;sfi;vn;s
\end{minipage}

\begin{minipage}[t]{0.45\textwidth}
    \centering
    \includegraphics[width=\linewidth]{dsc_cristallisation_chaude.pdf}
        \captionof{figure}{donees brutes}
        \label{fig:DSCprimitif}
\end{minipage}
\hfill
\begin{minipage}[h]{0.4\textwidth}
BLBBLBAAAA AAAAAAA AA AAAAAAAAAA,
skfvnlsnvf,sljvj
;sfi;vn;s
\end{minipage}

\begin{minipage}[t]{0.45\textwidth}
    \centering
    \includegraphics[width=\linewidth]{dsc_cristallisation_froide.pdf}
        \captionof{figure}{donees brutes}
        \label{fig:DSCprimitif}
\end{minipage}
\hfill
\begin{minipage}[h]{0.4\textwidth}
BLBBLBAAAA AAAAAAA AA AAAAAAAAAA,
skfvnlsnvf,sljvj
;sfi;vn;s
\end{minipage}

